% results_related_work

\section{Related Work}
\label{sec:related-work}

We draw on four strands: religious and devotional technology in HCI;
gamification and self-tracking, and specifically the literature on
tracking-induced guilt and anxiety that RQ1 contradicts; app-store review
mining as a research method; and the literature on feature bloat and
interface complexity that RQ5 tests.

\subsection{Religious and Devotional Technology in HCI}

HCI and CSCW have taken up religious practice as a first-class design concern
across several traditions and platforms. \citet{abokhodair2020holytweets}
analyse 2.6 million tweets sharing Quran verses on Arabic Twitter, combining
quantitative topic analysis with a qualitative study of users' devotional
goals, and demonstrate that large-scale, mixed-methods analysis of Islamic
digital practice is both feasible and productive --- a precedent our
ecosystem-scale review corpus follows for a different platform.
\citet{wolf2023blessing} take a research-through-design approach to a
Christian blessing ritual, designing for ``uncontrollability'' as a value in
tension with the on-demand interaction paradigms most software assumes.
\citet{bhattacharjee2025residual} interview Hindu migrants in Canada about
how they use social media and videoconferencing to sustain religious practice
and community across displacement. None of these studies addresses prayer-time
or devotional utility apps specifically, or reviews at ecosystem scale.

Closer to our object of study, \citet{bellar2017ipray} conducts a dissertation
study of Catholic and Islamic mobile prayer applications, textually analysing
65 app descriptions and user-testing one app from each tradition to identify
design approaches and the centrality of reminder features --- a finding our
RQ3 and RQ6 corroborate independently, at a different scale and with a
different method. \citet{azahari2025comparative} compare two named Islamic
applications, Muslim Pro and Sajda, through expert heuristic evaluation,
finding intrusive advertising and complex navigation in the feature-rich app
against interface simplicity and limited functionality in the minimal one ---
a pattern our RQ5 result revisits with corpus evidence rather than expert
judgement, and complicates: we do not find that a fuller feature set predicts
worse sentiment.

The closest methodological precedent to this paper is a pair of studies by
\citet{noor2025halaltourism} and \citet{noor2026halalsuperapps}, who apply
Leximancer content analysis to user reviews of a single Muslim-lifestyle
mobile application, identifying service-quality themes that include prayer
time accuracy, adhan reliability, and advertising --- themes that recur
throughout our own results. Both studies analyse one application (279 and
2,459 reviews respectively); ours analyses 26 applications and 732,194
reviews with aspect-based sentiment, seeded topic modelling, demand-signal
extraction, and mixed-effects statistical models rather than content-analysis
themes alone, which is what lets us distinguish an ecosystem-level finding
from an artefact of any one application's release history
(\S\ref{sec:rq1-coding}). We are not aware of a published ecosystem-scale,
mixed-methods study of devotional-app reviews at our corpus's scale, and we
position our contribution as filling that specific gap rather than as the
first study of Islamic mobile applications in HCI.

\subsection{Gamification and Self-Tracking}

A substantial literature documents that gamified self-tracking --- streaks,
points, badges --- can produce exactly the harm RQ1 tests for. In a CHI study
of habit-formation apps, \citet{renfree2016dontkick} show that streaks and
reminders are effective at sustaining new behaviour but create a dependency
whose loss introduces fragility into a user's change effort.
\citet{weathers2023activitystreaks} characterise what makes an activity
streak psychologically sticky, including the effect of framing a streak as a
daily rather than cumulative task. In interview studies of wearable-tracker
users, both \citet{humphreys2026toxic} and
\citet{demetrovics2025lifestylepressures} report internal conflict organised
around anxiety, guilt, loss of bodily autonomy, and what the latter study
names ``behavioral conformity and goal obsession.'' Experimentally,
\citet{bai2024darkside} find that competition-based gamification in fitness
apps increases stress relative to a self-exploration-oriented design. Against
this, a meta-analysis by \citet{six2021gamification} finds no significant
difference in depressive-symptom outcomes between gamified and non-gamified
mental-health apps, showing the harm is not universal even within
self-tracking generally.

We engage this literature on its strongest terms: the mechanisms it
identifies (dependency on an artificial streak, anxiety at a data mismatch
with lived experience, conformity to an externally imposed goal) are
well-evidenced in fitness and general habit-tracking contexts. RQ1 asks
whether the same mechanism operates in a devotional prayer tracker, where the
tracked behaviour is a religious obligation rather than a personal fitness
goal. We find that it largely does not: guilt language is under-represented
rather than over-represented in tracker discourse, and where a genuinely
normative complaint surfaces, its content is different in kind from anxiety
about a broken streak (\S\ref{sec:rq1}). We read this as a boundary condition
on the self-tracking guilt literature rather than as a refutation of it: the
harm this literature documents may depend on the tracked behaviour being
one the user set for themselves, which does not describe a prayer obligation.

\subsection{App-Store Review Mining as Method}

Mining app-store reviews for software-engineering and HCI insight is an
established method with its own maturing literature.
\citet{dabrowski2022analysing} survey 182 papers on app review analysis
published between 2012 and 2020, classifying the field by the information
mined, the techniques applied, and the software-engineering activity
supported; \citet{alsubaihin2019appstore} study how the app store as a
channel reshapes developers' own release and requirements practices;
\citet{ciurumelea2017analyzing} build a review taxonomy validated across 39
mobile applications; and \citet{memon2023appstoremining} combine topic
modelling, sentiment analysis, and functional/non-functional classification
in a pipeline structurally similar to our own, applied to a general
consumer-app sample.

The closest precedent for our scale and aspect-level design is
\citet{haggag2022mhealth}, who extract over five million reviews for 278
mHealth applications, classify them into 14 aspect categories, and validate by
hand-analysing 1,000 reviews per aspect. We follow their basic move ---
ecosystem-scale mining resolved to named aspects rather than global sentiment
--- in a domain their sample does not cover, and depart from them on one point
that matters for RQ2: several of their headline findings are stated as shares
of users awarding one star to an aspect, which treats the star rating as a
readable index of severity. Our corpus does not support that assumption
(\S\ref{sec:rq2}). We read this as a boundary on rating-indexed severity rather
than a correction, since the effect may be specific to defects users are
unwilling to let tank an app they otherwise value. Their gold samples are also
far larger than ours, which is why we report per-aspect precision with
intervals (\S\ref{sec:method-measurement}) rather than a single accuracy
figure.

\citet{guluzade2024eating} pair a functionality audit of 27 mindfulness-eating
and eating-disorder applications with a content analysis of 1,248 reviews, in a
domain where the tracked behaviour is personally sensitive. Their two-source
design, a hand-built functionality matrix read against what reviewers say, is
structurally our promise-versus-delivery comparison
(\S\ref{sec:method-promise}), and their concern with tracking that
misrepresents a user's situation is the nearest analogue we have found to RQ4's
menstrual-exemption gap outside a religious setting. Our contribution relative
to theirs is scale and inference, which is what lets us report how weakly the
two promise sources actually agree rather than treat a functionality matrix as
ground truth.

Our pipeline sits within this tradition, and this paper
contributes to it in two specific ways beyond the devotional domain: a
demonstration that star ratings systematically hide a class of substantive
complaint (\S\ref{sec:rq2}), and a demonstration that single-app qualitative
sampling can manufacture themes that read as genre-level findings
(\S\ref{sec:rq1-coding}). Both are risks inherent to the method that this
literature has not, to our knowledge, quantified directly.

\subsection{Feature Bloat and Interface Complexity}

That more features degrade user experience is a widely held assumption in
software and product design, more often asserted than tested.
\citet{marzi2022overfeaturing} reviews the ``over-featuring'' literature
across the terms it goes by --- feature creep, feature fatigue, overdesign,
scope creep, gold-plating --- and finds the phenomenon ranked among the top
ten risks of new-product-development failure alongside ``a lack of
theoretical development, and a limited investigation of [its] causes,
drivers, and performance effects.'' In one of the few empirical
investigations of the claim, \citet{blackler2016rtfm} find that
over-featured consumer products do produce negative emotional experiences
and that most users do not exercise most of a product's features. RQ5 tests
the specific, commonly assumed causal link --- that a higher feature count
predicts more bloat-related complaints --- directly against corpus data, and
finds it fails in both directions: feature count does not predict bloat
complaints, and the relationship that does hold runs opposite to the
assumption (\S\ref{sec:rq5}). Read alongside Marzi's observation that the
phenomenon is asserted more often than investigated, our result suggests the
assumption has outrun the evidence for it, at least in this domain.
